En conclusión, los diferentes tipos de algoritmos utilizados pueden tener complejidades teóricamente parecidas como $std::sort$, $MergeSort$ o $QuickSort$, pero en la práctica podemos notar que existen variaciones temporales (figura 1) o espaciales (figura 6) que nos indican que el análisis de la situación es importante para no elegir algoritmos subóptimos.
Un ejemplo de esto se observa en las figuras 2 y 3, donde $PatienceSort$ pasa de ser malo en arreglos ascendentes a óptimo en arreglos descendentes.
Por otro lado, el mayor contraste entre la teoría y la realidad se puede observar con los algoritmos de multiplicación de matrices, donde $Strassen$ parece temporalmente superior a $Naive$, en cambio la práctica nos demuestra lo contrario (figuras 7, 8, 9, 10 y 11) para los casos de prueba utilizados y nos esclarece las consecuencias que tiene ignorar conceptos que en la teoría se omiten.
